\chapter{Données}

Dans le cadre de ce mémoire, les données textuelles exploitées reposent sur trois ressources principales~:
le projet ParCoLab, l’Atlas sonore des langues régionales du LIMSI, et une traduction de la parabole du
fils prodigue. Ces corpus ont été choisis pour leur richesse linguistique, correspondant aux attentes de
notre projet. En effet, ces 3 projets s'intéressent aux langues régionales de France, particulièrement l'alsacien,
le corse et l'occitan.

Le corpus ParCoLab (ParCoLab – Multilingual Parallel Corpus) \citep{stosic_2024} constitue la base la plus
volumineuse. ParCoLab est un projet parallèle multilingue porté par le laboratoire CLLE
(Université de Toulouse) et soutenu par la Délégation générale à la langue
française et aux langues de France (DGLFLF). Celui-ci vise à
constituer et diffuser un corpus de textes alignés dans différentes
langues, dont des langues régionales de France centrales du projet.
Ces textes sont alignés au niveau des phrases, ce qui permet d’effectuer des analyses comparatives fines.

L’Atlas sonore des langues régionales du LIMSI (Laboratoire LIMSI, CNRS) \citep{boulademareuil_2017} est
une ressource patrimoniale remarquable. Ce projet consiste en une
carte interactive de France (voir figure \ref{atlas_limsi_fig}) sur laquelle sont disséminées plus de 200 communes où
des locuteurs ont enregistré des versions d’un même texte~: la fable d’Ésope \textit{La bise et le soleil}.
Ces enregistrements, souvent retranscrits dans des graphies adaptées, permettent d’étudier des dialectes
des diverses langues régionales françaises, et d’analyser la variation phonétique, lexicale ou morphologique
dans un contexte de terrain. Grâce à cette ressource, nous avons pu intégrer des textes en occitan très variés
phonétiquement, enrichissant le corpus avec des variations dialectales.

\begin{figure}[H]
    \centering
    \includegraphics[width=0.75\textwidth]{atlas_limsi.png}
    \caption{Carte de France de l'atlas LIMSI provenant du \href{https://atlas.lisn.upsaclay.fr/}{site de L'atlas}.}
    \label{atlas_limsi_fig}
\end{figure}

Enfin, nous utilisons une traduction de la 
\href{https://gallica.bnf.fr/ark:/12148/bd6t5334045p/f1.item#}{parabole du fils prodigue}\footnote{Traductions provenant du livre l'enfant prodigue en 88 patois divers de France, par Léopold FAVRE}
dans divers dialectes régionaux.
Cette parabole est un classique des enquêtes dialectologiques~: elle a été traduite dans de nombreuses 
langues régionales, notamment via les enquêtes menées par Coquebert de Montbret au XIX\up{e} siècle. 
Sur le plan linguistique, la version occitane de cette parabole (ainsi que des variantes en corse ou d’autres langues) 
offre un point d’ancrage commun (texte pivot) qui facilite la comparaison entre dialectes et la détection de
structures grammaticales, lexicales ou morphologiques typiques de chaque variété.

\section{Analyse de données}

\subsection{Données parallèles}

\begin{table}[H]
\centering
\begin{tabular}{lSSSSS}
\hline
\textbf{Données} & \textbf{Atlas LIMSI} & \textbf{Parcolab} & \textbf{Parabole du fils prodigue} & \textbf{Total} \\
\hline
Tokens & 7\ 922 & 35\ 371 & 11\ 729 & 55\ 022 \\
Phrases & 305 & 1\ 485 & 537 & 2\ 327 \\
Lemmes fr & 57 & 2\ 791 & 191 & 3\ 039 \\
Textes fr & 1 & 9 & 1 & 11 \\
Textes co & 6 & 9 & \NA & 15 \\
Textes gsw & 7 & 9 & \NA & 16 \\
Textes oc & 47 & 9 & 21 & 77  \\
\hline
\end{tabular}
\\
\smallskip
(fr: français, co: corse, gsw: alsacien, oc: occitan)
\smallskip
\caption{Données statistiques des trois corpus parallèles.}
\label{table_donnees_para}
\end{table}

Les données du tableau \ref{table_donnees_para} sont issues d'un programme Python fait par nos soins.
Ce dernier utilise la bibliothèque Spacy de manière à trouver le nombre (toutes langues confondues)
de phrases, de tokens (hors ponctuations, avec apostrophes et contractions séparées en deux tokens),
ainsi que le nombre de lemmes français. Aucun modèle Spacy n'existe pour l'alsacien, le corse et l'occitan,
ce qui rend impossible une lemmatisation. Cette difficulté est exacerbée, car ces langues sont principalement orales
et ne possèdent pas de norme orthographique stable.
On retrouve dans notre corpus une majorité de textes en occitan (77), mais aussi 15 textes en corse et 16
textes en alsacien. Cette répartition déséquilibrée s'explique principalement par la richesse dialectale de
l'occitan, dont les nombreuses variétés régionales sont largement documentées dans l'Atlas LIMSI
(dans notre cas surtout en provenance du Languedoc et de Nouvelle-Occitanie).
Le fait que ParCoLab soit le corpus le plus textuellement riche (dû à la taille des textes
présents dans celui-ci) explique qu’il présente également le nombre de tokens le plus élevé.

D'un point de vue qualitatif, les trois corpus présentent des profils linguistiques très distincts qui se
complètent utilement. ParCoLab offre des textes alignés au niveau des phrases dans un cadre méthodologique
rigoureux, ce qui en fait la ressource la plus fiable pour l'entraînement du modèle. L'Atlas LIMSI,
quant à lui, apporte une richesse dialectale considérable, notamment pour l'occitan dont les 47 variétés
couvrent un spectre phonétique et lexical très large. Les transcriptions dialectales pouvant présenter
des graphies très hétérogènes d'un locuteur à l'autre, cette variété amène aussi plus de bruit dans les données
d'entraînement. La parabole du fils prodigue, elle, joue le rôle de texte pivot.
En effet, sa structure narrative fixe facilite la comparaison inter-dialectale, mais limite en
contrepartie la diversité lexicale et thématique des données.

\begin{table}[H]
\centering
\begin{tabular}{lSSSS}
\hline
\textbf{Données} & \textbf{Français} & \textbf{Alsacien} & \textbf{Corse} & \textbf{Occitan} \\
\hline
Fichiers & 11 & 16 & 15 & 77 \\
Tokens & 10\ 070 & 9\ 677 & 9\ 601 & 25\ 674 \\
Phrases & 363 & 423 & 418 & 1\ 123 \\
Lemmes & 3\ 039 & \NA & \NA & \NA \\
\hline
\end{tabular}
\caption{Données statistiques du corpus parallèles par langues.}
\label{table_donnees_par_langues}
\end{table}

Le tableau \ref{table_donnees_par_langues} permet d'affiner cette analyse en comparant la répartition des
données par langue. On constate que l'occitan domine très largement avec 25\,674 tokens répartis sur 77 fichiers,
contre respectivement 9\,677 et 9\,601 tokens pour l'alsacien et le corse. Cette surreprésentation de
l'occitan peut nous poser problème. Celle-ci pourrait conduire à un classifieur moins performant sur les
paires alsaciennes ou corses, qui disposent de moins d'exemples d'entraînement.
Le corse et l'alsacien sont très similaires statistiquement. Cette similarité nous indique que les deux
langues seront couvertes de manière quasi-égalitaire. Enfin, le nombre de lemmes français (3\,039) constitue
un indicateur indirect de la richesse lexicale des textes sources. En théorie, ce chiffre devrait être
similaire quelle que soit la langue cible, les textes étant des traductions d'un même contenu.
Cependant la lemmatisation des autres langues reste pour l'instant impossible à cause du manque de modèle Spacy
pour celles-ci.

\subsection{Données comparables}

\begin{table}[H]
\centering
\begin{tabular}{lSSSS}
\hline
\textbf{Données} & \textbf{Alsacien} & \textbf{Corse} & \textbf{Francais} & \textbf{Occitan} \\
\hline
Fichiers & 1\ 590 & 5 & 2\ 760 & 1\ 590 \\
Tokens & 269\ 277 & 352 & 6\ 205\ 620 & 311\ 734 \\
Phrases & 16\ 940 & 19 & 278\ 025 & 14\ 198 \\
Lemmes & \NA & \NA & 1\ 868\ 027 & \NA \\
\hline
\end{tabular}
\caption{Données statistiques du corpus comparable Wikipédia.}
\label{table_donnees_comp}
\end{table}

Le tableau \ref{table_donnees_comp} représente les données du corpus comparable que nous utiliserons pour la
phase de validation du modèle que nous allons créer.
Celles-ci ont été obtenues via Wikipédia et les liens interlangues.
Le programme va rechercher les articles en français ayant un lien interlangue avec l'alsacien et/ou le corse et/ou l'occitan.
Dans le cas de l'alsacien, le programme cherche uniquement les articles en alsacien et non dans les autres langues allémaniques en s'aidant
des liens Wikipédia, ceux en alsacien provenant de l'adresse \href{https://als.wikipedia.org/}{https://als.wikipedia.org/}.
Pour ce faire, nous avons utilisé le programme\footnote{\href{https://gitlab.com/memoire5005040/fragments_paralleles}{Le programme "fragment\_paralleles" de P. Tsolakis}} de \citet{tsolakis_2024}.
Le programme original a pour but de récupérer
les articles en corse, luxembourgeois et occitan ainsi qu'en français, anglais, allemand, italien et espagnol.
Ensuite le programme sélectionne 10 phrases provenant de 15 articles aléatoires pour chaque langue traitée
de manière à trouver des paires de phrases candidates.
Nous n'utilisons que la première partie du programme (modifiée pour les langues que nous étudions) car elle nous permet de retrouver
tous les articles ayant un lien interlangue entre le français et les langues régionales~: alsacien, corse et occitan.

Le français domine massivement avec plus de 6 millions de tokens répartis sur 2\,760 fichiers, reflétant la
richesse et la maturité de la communauté Wikipédia francophone. Cette domination du français vient du fait que certaines
des pages ne sont pas traduites en alsacien mais sont présentes en corse et/ou occitan, le fichier français est donc ajouté
au corpus, ainsi que ses traductions dans les langues régionales (hors alsacien). Ces fichiers étant présents, le français devient
la langue majoritaire car l'on ne cherche pas que les articles où les 3 langues sont présentes en même temps.
L'alsacien, avec 269\,277 tokens sur 1\,590
fichiers, se démarque nettement des deux autres langues régionales et témoigne de la contribution
de ses locuteurs, qui intègre notamment les dialectes mulhousien et strasbourgeois (voir en quelle proportion dans le tableau \ref{table_dialectes_gsw}).

\begin{table}[H]
\centering
\begin{tabular}{lSSS}
\hline
\textbf{Données} & \textbf{Alsacien} & \textbf{Mulhousien} & \textbf{Strasbourgeois} \\
\hline
Fichiers & 1\ 194 & 392 & 4 \\
Tokens & 146\ 675 & 120\ 717 & 1\ 885 \\
Phrases & 10\ 860 & 5\ 993 & 87 \\
\hline
\end{tabular}
\caption{Tableau de décomposition de l'alsacien et de ses dialectes.}
\label{table_dialectes_gsw}
\end{table}

Cette variété dialectale interne à l'alsacien constitue un atout pour évaluer la robustesse du modèle face
à des graphies divergentes, mais introduit également une complexité supplémentaire que les données
d'entraînement parallèles, plus homogènes, ne couvrent qu'en moindre mesure.

L'occitan, avec 311\,734 tokens sur 1\,590 fichiers, présente un volume intermédiaire cohérent avec sa présence
sur Wikipédia.
Le corse, en revanche, apparaît très faiblement représenté avec seulement 352 tokens sur 5 fichiers,
soit un volume environ 60 fois inférieur à celui de l'alsacien. Ce volume extrêmement bas
d'articles doit être dû à un manque d'utilisateurs ajoutant des pages en corse sur Wikipédia. Cette disproportion est préoccupante pour
la phase de validation~: le modèle disposera de très peu d'exemples pour être évalué, ce qui limitera la
portée des conclusions tirées sur les performances du classifieur pour cette langue.

D'un point de vue qualitatif, les textes Wikipédia présentent une hétérogénéité thématique et stylistique
bien plus grande que les corpus parallèles. Cette diversité est à la fois un avantage et une difficulté.
Elle permet de tester la généralisation du modèle au-delà des genres textuels vus à l'entraînement, mais
elle expose aussi le classifieur à des registres, des vocabulaires spécialisés et des structures discursives
pour lesquels il n'a pas été préparé. Cette tension entre la richesse du corpus de validation et les
limites du corpus d'entraînement constitue une contrainte importante dont il faudra tenir
compte dans l'interprétation des résultats.